\section{Introdução}
\label{sec:introduction}

Problemas de otimização combinatória surgem quando é necessário escolher,
entre alternativas discretas, uma solução de menor custo que satisfaça
determinadas restrições. A quantidade de alternativas pode crescer
exponencialmente com a entrada. Nesses casos, a escolha de uma estratégia
algorítmica não envolve apenas obter uma solução: é preciso considerar
o tempo disponível, a qualidade exigida e a possibilidade de demonstrar
que não existe solução melhor \cite{cormen2022,fomin2010}.

O Problema da Cobertura Mínima por Vértices, ou \emph{Minimum Vertex Cover},
expressa essa tensão de forma particularmente clara. Dado um grafo, deseja-se
selecionar o menor conjunto de vértices que contenha pelo menos uma
extremidade de cada aresta. Essa condição permite modelar, sob hipóteses
apropriadas, a seleção de pontos para observar conexões, inspecionar
ligações de uma infraestrutura ou instrumentar relações entre componentes.
A versão de decisão é NP-completa, enquanto a versão de otimização é
NP-difícil \cite{karp1972,garey1979}.

A dificuldade no pior caso não impede estudar algoritmos exatos úteis em
instâncias limitadas. A técnica Branch and Bound organiza uma busca e
elimina ramos quando limites demonstram que eles não podem melhorar a
solução conhecida. Em contraste, algoritmos aproximados podem obter
rapidamente uma cobertura válida com uma garantia matemática sobre sua
cardinalidade. O método por emparelhamento maximal oferece garantia de
fator~2 \cite{cormen2022,vazirani2001}. São propostas com compromissos
distintos, cuja comparação exige não confundir uma boa solução com uma
prova de otimalidade.

O objetivo geral deste trabalho é analisar experimentalmente um Branch
and Bound e uma aproximação por emparelhamento maximal, implementados
em Java~17 puro, comparando tempo, qualidade das coberturas e esforço de
busca em grafos de diferentes tamanhos e densidades. Os objetivos
específicos são formalizar o problema e sua classificação; descrever e
validar os dois algoritmos; gerar instâncias reproduzíveis; medir seu
comportamento sob um limite de tempo; e discutir as vantagens e
limitações sustentadas pelos dados.

A Seção~\ref{sec:theory} apresenta os fundamentos, a demonstração de
NP-Completude e outras abordagens. A Seção~\ref{sec:methodology} descreve
a implementação, os testes e o protocolo. As Seções~\ref{sec:results}
e~\ref{sec:analysis} apresentam e analisam os resultados, respectivamente.
A Seção~\ref{sec:conclusion} sintetiza as conclusões. Os anexos incluem
o código-fonte essencial e instruções de reprodução.
